\chapter{Wnioskowanie z procesów zależnych: autokorelacja i modele szeregów czasowych}
\label{ch:time-series}
\paragraph{Powiązania i odsyłacze.} Podstawy procesów Markowa i stacjonarności przedstawia tom II (rozdziały~\ref{II-ch:ctmc} i~\ref{II-ch:gauss}); klasyczne modele regresji — rozdział~\ref{ch:linear-regression}.


\section{Dlaczego sto dni to nie zawsze sto niezależnych obserwacji?}
Pomiar temperatury w kolejnych godzinach, dzienne straty
finansowe lub odczyty jednego urządzenia tworzą \emph{szereg czasowy}:
uporządkowane obserwacje jednego procesu. Zamiana ich kolejności
zazwyczaj zmienia sens modelu. Jeżeli zależność w czasie jest
pomijana, nawet poprawna średnia z próby może mieć źle oszacowaną
wariancję, a losowa walidacja krzyżowa może wykorzystywać
informację z przyszłości. Podstawowe pojęcia procesów i
filtracji wprowadziliśmy w tomie II; obecnie naszym zadaniem
jest \emph{estymacja i predykcja} na podstawie jednej trajektorii.

\begin{definicja}[Słaba stacjonarność i autokowariancja]
\label{def:weak-stationarity}
Niech $(X_t)_{t\in\mathbb Z}$ będzie procesem rzeczywistym
o skończonych drugich momentach. Nazywamy go \emph{słabo
stacjonarnym}, jeśli istnieje $\mu\in\R$ takie, że
$\E X_t=\mu$ dla wszystkich $t$, a
$\Cov(X_t,X_s)=\gamma(t-s)$ dla pewnej funkcji
$\gamma:\mathbb Z\to\R$. Wtedy $\gamma(h)$ to
\emph{autokowariancja} opóźnienia $h$, a przy
$\gamma(0)>0$ wielkość $\rho(h)=\gamma(h)/\gamma(0)$
to \emph{autokorelacja}. Zachodzą $\gamma(-h)=\gamma(h)$
i $|\rho(h)|\le1$. Proces ściśle stacjonarny ma
niezmiennicze względem przesunięcia czasu rozkłady
skończeniewymiarowe; przy skończonych drugich momentach
jest też słabo stacjonarny.
\end{definicja}

\begin{uwaga}[Stacjonarność a trend]
Jeśli średnia rośnie w czasie, dane nie są słabo
stacjonarne, choć reszty po właściwym usunięciu trendu
mogą być stacjonarne. Również szeregi sezonowe wymagają
jawnego modelowania kalendarza. Automatyczne różnicowanie
każdej serii bez diagnozy może usunąć użyteczny sygnał.
\end{uwaga}

\begin{twierdzenie}[Dokładna wariancja średniej z szeregu]
\label{thm:mean-time-variance}
Niech $(X_t)_{t\in\mathbb Z}$ będzie słabo stacjonarny
o autokowariancji $\gamma$. Dla
$\overline X_n=n^{-1}\sum_{t=1}^nX_t$ zachodzi
\[
 \Var(\overline X_n)=\frac1{n^2}
 \left(n\gamma(0)+2\sum_{h=1}^{n-1}(n-h)\gamma(h)\right).
\]
Jeśli dodatkowo $\sum_{h\in\mathbb Z}|\gamma(h)|<\infty$,
to
$\lim_{n\to\infty}n\Var(\overline X_n)
=\sum_{h\in\mathbb Z}\gamma(h)=:\tau^2\ge0$.
\end{twierdzenie}
\begin{proof}
Z definicji wariancji sumy
$n^2\Var(\overline X_n)=\sum_{s=1}^n\sum_{t=1}^n
\Cov(X_s,X_t)$. Dla $s=t$ występuje $n$ składników
$\gamma(0)$; dla każdego $h\in\{1,\ldots,n-1\}$
istnieje $n-h$ par o $t-s=h$ i tyle samo o $s-t=h$.
Daje to pierwszy wzór. Mnożąc go przez $1/n$,
otrzymujemy
$n\Var(\overline X_n)=\sum_{|h|<n}
(1-|h|/n)\gamma(h)$. Bezwzględna sumowalność
pozwala zastosować zbieżność zdominowaną dla szeregów.
Nieujemność granicy wynika z nieujemności wariancji.
\end{proof}

\begin{przyklad}[Jak korelacja zmienia efektywną liczebność?]
Jeśli $\gamma(h)=\gamma(0)\phi^{|h|}$ dla $|\phi|<1$,
to $\tau^2=\gamma(0)(1+\phi)/(1-\phi)$.
Dla $\phi=0{,}8$ błąd standardowy średniej w dużej próbie
jest w przybliżeniu trzykrotnie większy niż przy błędnym
założeniu niezależności. Odpowiada to heurystycznej
liczebności efektywnej $n_{\mathrm{eff}}\approx n/9$.
Jest to \emph{porównanie wariancji}, nie dosłowne
stwierdzenie, że usunięto obserwacje. Przy ujemnej
autokorelacji liczebność efektywna może przekroczyć $n$.
\end{przyklad}

\section{Model AR(1): konstrukcja, estymacja i prognoza}
\begin{definicja}[Biały szum i model autoregresyjny]
Mówimy, że $(\varepsilon_t)_{t\in\mathbb Z}$ jest
\emph{białym szumem} o wariancji $\sigma^2>0$, jeśli
$\E\varepsilon_t=0$ oraz
$\Cov(\varepsilon_t,\varepsilon_s)=\sigma^2\ind{t=s}$.
W poniższym twierdzeniu założymy ponadto niezależność
oraz identyczny rozkład innowacji. \emph{AR(1)} opisuje
proces $(X_t)$ spełniający
$X_t-\mu=\phi(X_{t-1}-\mu)+\varepsilon_t$,
gdzie $\mu\in\R$ i $\phi\in\R$ są parametrami.
\end{definicja}

\begin{twierdzenie}[Stacjonarne rozwiązanie AR(1)]
\label{thm:ar1-stationary}
Niech $(\varepsilon_t)_{t\in\mathbb Z}$ będą
niezależnymi zmiennymi o wspólnym rozkładzie,
$\E\varepsilon_t=0$ i
$\Var(\varepsilon_t)=\sigma^2\in(0,\infty)$.
Niech $\mu\in\R$ i $|\phi|<1$. Wtedy szereg
$X_t=\mu+\sum_{j=0}^{\infty}\phi^j\varepsilon_{t-j}$
jest zbieżny w $L^2$ dla każdego $t$ i definiuje
słabo stacjonarne rozwiązanie AR(1) z
\[
 \E X_t=\mu,\qquad
 \gamma(h)=\frac{\sigma^2}{1-\phi^2}\phi^{|h|},
 \quad h\in\mathbb Z.
\]
Jest to jedyne słabo stacjonarne rozwiązanie o skończonym
drugim momencie na tej samej przestrzeni, napędzane tym samym
szumem i spełniające to równanie w $L^2$.
\end{twierdzenie}
\begin{proof}
Niezależność i centrowanie dają dla $m>k$
$\E|\sum_{j=k}^m\phi^j\varepsilon_{t-j}|^2
=\sigma^2\sum_{j=k}^m\phi^{2j}$.
Ponieważ $|\phi|<1$, sumy częściowe są ciągiem
Cauchy'ego w $L^2$. Przesunięcie indeksów pokazuje
$X_t-\mu=\phi(X_{t-1}-\mu)+\varepsilon_t$ w $L^2$.
Dla $h\ge0$ w iloczynie dwóch szeregów przeżywają
jedynie składniki ze wspólnym indeksem innowacji:
$\gamma(h)=\sigma^2\sum_{j\ge0}\phi^{2j+h}
=\sigma^2\phi^h/(1-\phi^2)$. Pozostałe opóźnienia
wynikają z symetrii kowariancji.
Jeżeli $X_t$ i $X'_t$ są dwoma słabo stacjonarnymi
rozwiązaniami z tymi samymi innowacjami, to
$D_t=X_t-X'_t=\phi^mD_{t-m}$ dla każdego $m$.
Słaba stacjonarność obu procesów zapewnia
$\sup_t\E D_t^2<\infty$ przez nierówność
$D_t^2\le2X_t^2+2(X'_t)^2$.
Stąd $\E D_t^2\le\phi^{2m}\sup_s\E D_s^2\to0$.
\end{proof}

\begin{uwaga}[Jaki jest sens parametru $\phi$?]
Dodatnie $\phi$ opisuje utrzymywanie się odchyleń
od poziomu $\mu$, a ujemne naprzemienność znaków.
Dla $\phi=0$ otrzymujemy niezależny szum wokół średniej.
Dla $\phi=1$ zwykły proces z sumowaniem innowacji
jest błądzeniem losowym, a nie stacjonarnym AR(1).
Przy $|\phi|\ge1$ nie istnieje opisane powyżej
przyczynowe stacjonarne rozwiązanie napędzane
jednostronną sumą przeszłych innowacji; nie wyklucza
to nietypowych rozwiązań z informacją o przyszłości.
\end{uwaga}

\begin{twierdzenie}[Prognoza AR(1) przy znanych parametrach]
\label{thm:ar1-forecast}
W założeniach twierdzenia~\ref{thm:ar1-stationary}
niech $\mathcal F_t=\sigma(\varepsilon_s:s\le t)$.
Dla każdego $h\in\N$, $h\ge1$, zachodzi
\[
 \E[X_{t+h}\mid\mathcal F_t]
 =\mu+\phi^h(X_t-\mu),\qquad
 \Var(X_{t+h}\mid\mathcal F_t)
 =\sigma^2\sum_{j=0}^{h-1}\phi^{2j}.
\]
Jeżeli innowacje są normalne, warunkowy rozkład jest
normalny i podane parametry wyznaczają dokładny
przedział predykcji przy \emph{znanych} $\mu,\phi,\sigma$.
\end{twierdzenie}
\begin{proof}
Iteracja rekurencji prowadzi do
$X_{t+h}-\mu=\phi^h(X_t-\mu)+
\sum_{j=0}^{h-1}\phi^j\varepsilon_{t+h-j}$.
Innowacje w ostatniej sumie są niezależne od
$\mathcal F_t$, zatem ich warunkowa średnia
wynosi zero, a wariancja to suma wariancji.
Przy normalności suma jest normalna.
\end{proof}

\begin{stwierdzenie}[Estymacja warunkowa AR(1)]
\label{prop:ar1-conditional-mle}
Załóżmy, że obserwujemy $x_1,\ldots,x_n$, $n\ge3$,
z modelu $X_t=\phi X_{t-1}+\varepsilon_t$ dla $t\ge2$,
gdzie $\varepsilon_t$ są niezależne
$\mathcal N(0,\sigma^2)$ i niezależne od $X_1$.
Przy znanej średniej zerowej i nieznanych
$\phi\in\R$, $\sigma^2>0$, dla każdej ustalonej
$\sigma^2$ log-wiarygodność warunkowa względem $X_1=x_1$
osiąga maksimum po $\phi$ w punkcie
\[
 \widehat\phi_{\mathrm{CLS}}=
 \frac{\sum_{t=2}^n x_tx_{t-1}}{\sum_{t=2}^n x_{t-1}^2},
\]
jeśli mianownik jest dodatni. Estymator ten jest
również estymatorem MNK dla regresji bez wyrazu wolnego.
Jeżeli dodatkowo resztowa suma kwadratów jest dodatnia,
wspólne maksimum po $\phi$ i $\sigma^2>0$ istnieje,
a warunkowy MLE wariancji to $\mathrm{RSS}/(n-1)$.
Jeżeli zaś $\mathrm{RSS}=0$, wiarygodność rośnie
bez ograniczenia przy $\sigma^2\downarrow0$ i wspólny MLE
w otwartej przestrzeni parametrów nie istnieje.
\end{stwierdzenie}
\begin{proof}
Warunkowa log-wiarygodność, z pominięciem stałej,
wynosi $-\frac{n-1}{2}\log v-
\frac{1}{2v}\sum_{t=2}^n(x_t-\phi x_{t-1})^2$.
Dla ustalonego $v$ jej maksimum po $\phi$ odpowiada
minimum wypukłej funkcji kwadratowej; równanie
normalne daje podany iloraz. Maksymalizacja po
$v>0$ odbywa się tak samo jak w twierdzeniu
\ref{thm:aic-gaussian}, z liczebnością $n-1$.
\end{proof}

\begin{uwaga}[Warianty białego szumu]
Biały szum w definicji powyżej wymaga jedynie nieskorelowania
innowacji. Z samego nieskorelowania nie wolno wyprowadzić
niezależności warunkowej, formuły wiarygodności Gaussowskiej
ani prognozy jako pełnego rozkładu warunkowego. W twierdzeniach
o stacjonarnym AR(1) i prognozie założyliśmy dlatego
jawnie niezależność innowacji; do warunkowej wiarygodności
potrzebna jest ponadto ich normalność.
\end{uwaga}

\begin{uwaga}[Warunkowe a dokładne MLE]
Gdy model wymaga $|\phi|<1$, powyższe rozwiązanie
bez ograniczenia może leżeć poza przestrzenią parametrów;
wtedy należy maksymalizować na dopuszczalnej dziedzinie
(w otwartym przedziale supremum może nie być osiągnięte).
Dokładna wiarygodność stacjonarnego normalnego AR(1)
uwzględnia ponadto rozkład $X_1$ i może dawać inny
estymator. Niepewność \emph{oszacowanych} parametrów
nie jest uwzględniona w przedziale z twierdzenia
\ref{thm:ar1-forecast}; w małych próbach ma to znaczenie.
\end{uwaga}

\section{MA, ARMA, ACF i identyfikacja struktury}
\begin{definicja}[Procesy MA($q$) i ARMA($p,q$)]
Niech $(\varepsilon_t)_{t\in\mathbb Z}$ będzie
białym szumem o wariancji $\sigma^2>0$.
Dla $q\in\N\cup\{0\}$ i współczynników
$\theta_1,\ldots,\theta_q\in\R$ proces
$X_t=\mu+\varepsilon_t+
\sum_{j=1}^q\theta_j\varepsilon_{t-j}$
to \emph{MA($q$)}. Wprowadzamy operator opóźnienia
$BX_t=X_{t-1}$ i wielomiany
$\Phi(z)=1-\sum_{i=1}^p\phi_i z^i$ oraz
$\Theta(z)=1+\sum_{j=1}^q\theta_j z^j$.
Model \emph{ARMA($p,q$)} (autoregressive moving average,
czyli połączenie autoregresji i średniej ruchomej) jest równaniem
$\Phi(B)(X_t-\mu)=\Theta(B)\varepsilon_t$.
\emph{Przyczynowość} oznacza możliwość wyrażenia
$X_t-\mu$ jako zbieżnego w $L^2$ szeregu
z przeszłych i bieżących innowacji; \emph{odwracalność}
oznacza analogiczną możliwość odzyskania innowacji
z przeszłych i bieżących wartości procesu.
\end{definicja}

\begin{twierdzenie}[Autokowariancja MA($q$)]
\label{thm:maq-cov}
Niech $X_t-\mu=\sum_{j=0}^q\theta_j\varepsilon_{t-j}$,
$\theta_0=1$, a innowacje będą nieskorelowane,
wycentrowane i mają wspólną wariancję $\sigma^2$.
Wtedy proces jest słabo stacjonarny i dla
$0\le h\le q$ zachodzi
$\gamma(h)=\sigma^2\sum_{j=0}^{q-h}\theta_j\theta_{j+h}$,
natomiast $\gamma(h)=0$ przy $h>q$.
\end{twierdzenie}
\begin{proof}
Rozwijamy
$\Cov(X_t,X_{t-h})=\sum_{j,k=0}^q
\theta_j\theta_k\Cov(\varepsilon_{t-j},\varepsilon_{t-h-k})$.
Dodatnie składniki otrzymujemy tylko, gdy
$j=h+k$, co daje $\sigma^2\sum_{k=0}^{q-h}
\theta_{h+k}\theta_k$. Gdy $h>q$, żadne indeksy
nie mogą być równe. Kowariancja nie zależy od $t$.
\end{proof}

\begin{uwaga}[Charakterystyki ACF i PACF]
\emph{ACF} (autocorrelation function) to funkcja autokorelacji
$\rho(h)$. \emph{PACF} (partial autocorrelation function),
czyli częściowa funkcja autokorelacji, na opóźnieniu $h$
to korelacja liniowych
reszt z regresji $X_t$ i $X_{t-h}$ na pośrednie
$X_{t-1},\ldots,X_{t-h+1}$, jeżeli kowariancja
wektora pośredniego jest odwracalna; dla $h=1$
PACF jest zwykłą autokorelacją opóźnienia $1$.
Dla MA($q$) teoretyczna ACF zeruje się po $q$,
zaś dla AR($p$) w standardowym modelu przyczynowym
teoretyczna PACF zeruje się po $p$. W skończonej
próbie \emph{oszacowane} wartości nie zerują się dokładnie,
więc nie jest to mechaniczna reguła identyfikacji rzędu.
\end{uwaga}

\begin{stwierdzenie}[Warunki przyczynowości i odwracalności]
Jeśli wielomian $\Phi$ modelu ARMA($p,q$) nie ma zer
w kole domkniętym $\{z\in\mathbb C:|z|\le1\}$,
to rozwinięcie $\Theta(z)/\Phi(z)=\sum_{j=0}^\infty\psi_jz^j$
ma współczynniki bezwzględnie sumowalne i definiuje
przyczynowe rozwiązanie
$X_t-\mu=\sum_{j\ge0}\psi_j\varepsilon_{t-j}$
w $L^2$, dla białego szumu o skończonej wariancji.
Jeśli dodatkowo $\Theta$ nie ma zer w tym kole,
a wielomiany $\Phi$ i $\Theta$ nie mają wspólnych czynników,
innowacje można odzyskać przez bezwzględnie sumowalny
filtr z obserwacji (odwracalność).
\end{stwierdzenie}
\begin{proof}
Jeśli wielomian nie zeruje się na kole o promieniu
większym od $1$, odwrotność $1/\Phi(z)$ jest
holomorficzna na tym kole. Szereg Taylora ma wtedy
współczynniki malejące co najmniej geometrycznie
(ze wzoru Cauchy'ego na pochodne), więc jest
bezwzględnie sumowalny. Iloczyn z wielomianem
$\Theta$ zachowuje tę własność. Dla białego szumu
odpowiadający szereg jest Cauchy'ego w $L^2$.
Mnożenie formalnych szeregów, uzasadnione
sumowalnością, odtwarza równanie ARMA.
Dla $\Theta$ stosujemy identyczny argument do
$\Phi/\Theta$, otrzymując filtr odwrotny.
\end{proof}

\section{Diagnostyka autokorelacji i niepewność estymacji}
Dla obserwacji $x_1,\ldots,x_n$ definiujemy
$\widehat\gamma(h)=n^{-1}\sum_{t=h+1}^n
(x_t-\overline x)(x_{t-h}-\overline x)$,
$0\le h<n$, a przy $\widehat\gamma(0)>0$
$\widehat\rho(h)=\widehat\gamma(h)/\widehat\gamma(0)$.
Zwyczajowo rysujemy ACF zarówno oryginalnych danych,
jak i reszt dopasowanego modelu. Użycie wspólnego
mianownika $n$ jest jedną z konwencji; oszacowania ACF
mogą zależeć od wyboru normalizacji i średniej.

\begin{definicja}[Statystyka Ljunga--Boxa]
Dla szeregu długości $n$ i opóźnień $1\le m<n$
przy $\widehat\gamma(0)>0$ statystykę portmanteau
definiujemy wzorem
\[
 Q_m=n(n+2)\sum_{h=1}^{m}\frac{\widehat\rho(h)^2}{n-h}.
\]
Stosujemy ją do badania \emph{łącznej} obecności
autokorelacji do opóźnienia $m$. W klasycznym
przybliżeniu dla reszt dopasowanego ARMA($p,q$)
porównuje się ją z rozkładem $\chi^2_{m-p-q}$,
o ile $m>p+q$ i spełnione są dodatkowe warunki
asymptotyczne modelu i procedury estymacji.
\end{definicja}

\begin{uwaga}[Jak interpretować test Ljunga--Boxa?]
Duża wartość p nie dowodzi braku autokorelacji;
przy małej próbie test ma ograniczoną moc.
Z kolei istotność może wynikać ze zmiennej
wariancji, niestacjonarności lub błędnej specyfikacji.
Wybór $m$ należy uzasadnić częstotliwością danych,
sezonowością i horyzontem prognozy, nie zaś dokonać
go po wielokrotnym przeszukiwaniu wartości p.
\end{uwaga}

\begin{stwierdzenie}[Korekta HAC: sens estymatora wariancji długookresowej]
\label{prop:hac}
Niech $(X_t)$ będzie słabo stacjonarny o bezwzględnie
sumowalnej autokowariancji. Wówczas wariancja
asymptotyczna $\sqrt n(\overline X_n-\mu)$ wynosi
$\tau^2=\sum_{h\in\mathbb Z}\gamma(h)$, w sensie
granicy wariancji z twierdzenia~\ref{thm:mean-time-variance}.
Typową estymatą tej wielkości jest jądrowy estymator
Bartletta--Neweya--Westa; określenie HAC oznacza
\emph{heteroskedasticity and autocorrelation consistent},
czyli odporność na heteroskedastyczność i autokorelację
przy odpowiednich warunkach asymptotycznych:
\[
 \widehat\tau^2_b=\widehat\gamma(0)+
 2\sum_{h=1}^{b}\left(1-\frac{h}{b+1}\right)
 \widehat\gamma(h),\qquad 0\le b<n.
\]
Jego \emph{zgodność} nie wynika z samej stacjonarności:
wymaga dodatkowych warunków zależności, momentów oraz
odpowiednio wolnego wzrostu $b$ z liczebnością próby.
\end{stwierdzenie}

\begin{uwaga}[Wariancja nie jest centralnym twierdzeniem granicznym]
Z równości $n\Var(\overline X_n)\to\tau^2$
nie wynika sama z siebie asymptotyczna normalność.
Potrzebne jest osobne centralne twierdzenie graniczne
dla klasy rozważanych procesów, np. poprawnie określonej
klasy procesów mieszających. Bez niego przedział
$\overline X_n\pm z_{1-\alpha/2}\widehat\tau_b/\sqrt n$
nie ma automatycznej gwarancji pokrycia. Dla
normalnego AR(1) normalność średniej jest dokładna,
ale szacowanie $\tau$ pozostaje dodatkowym problemem.
\end{uwaga}

\section{Praktyczny wybór modelu i uczciwa walidacja w czasie}
Nie dzielimy danych chronologicznych poprzez losowe
przemieszanie, jeśli celem jest prognoza przyszłości.
Jeśli $\Delta X_t=X_t-X_{t-1}$, to model
ARIMA($p,d,q$) zakłada, że $\Delta^d X_t$ spełnia
odpowiedni model ARMA($p,q$); $d$ jest liczbą
różnicowań, a w odmianie sezonowej uwzględnia się
dodatkowo opóźnienia sezonowe. Takiego różnicowania
nie należy wykonywać bez sprawdzenia potrzeby i
przyjętej interpretacji procesu.
Na odłożonym bloku przyszłych obserwacji często mierzymy
RMSE (root mean squared error), czyli pierwiastek
średniego kwadratu błędu prognozy:
$\operatorname{RMSE}=\sqrt{m^{-1}\sum_{j=1}^m(y_j-\widehat y_j)^2}$,
gdzie $m$ jest liczbą sprawdzanych prognoz.
W \emph{walidacji przesuwającego się początku} uczymy
na przeszłości i prognozujemy późniejszy odcinek.
Dwa typowe warianty to poszerzające się okno treningowe
oraz przesuwające się okno o stałej długości.
To drugie może lepiej adaptować się do zmian
warunków, ale korzysta z mniejszej ilości danych.

\begin{center}
\begin{tikzpicture}[x=0.85cm,y=0.53cm,every node/.style={font=\small}]
\foreach \r/\k in {0/4,1/6,2/8} {
  \draw[fill=blue!13,draw=blue!50!black] (0,-\r*1.55) rectangle (\k,-\r*1.55+0.85);
  \draw[fill=orange!22,draw=orange!60!black] (\k,-\r*1.55) rectangle (\k+2,-\r*1.55+0.85);
}
\node[anchor=west] at (8.3,0.48) {trening};
\node[anchor=west] at (8.3,0.02) {prognoza};
\draw[fill=blue!13,draw=blue!50!black] (7.65,0.57) rectangle (8.0,0.79);
\draw[fill=orange!22,draw=orange!60!black] (7.65,0.10) rectangle (8.0,0.32);
\node[anchor=west] at (0,-4.5) {Każde okno testowe następuje po swoim oknie treningowym.};
\end{tikzpicture}
\end{center}

\begin{lstlisting}[language=Python,caption={Kompletny program: symulacja AR(1), estymacja, test reszt i prognozy chronologiczne (Python).}]
import numpy as np
from scipy.stats import norm
from statsmodels.stats.diagnostic import acorr_ljungbox

rng = np.random.default_rng(2026)
n = 260
phi_true, sigma_true = 0.75, 1.0
noise = rng.normal(0.0, sigma_true, size=n + 300)
series = np.zeros(n + 300)
for t in range(1, len(series)):
    series[t] = phi_true * series[t-1] + noise[t]
x = series[300:]  # burn-in: reduces dependence on x[0] = 0

# Chronological holdout (last 40 points) and one-step forecast.
train = x[:220]
test = x[220:]
phi_hat = np.dot(train[1:], train[:-1]) / np.dot(train[:-1], train[:-1])
resid = train[1:] - phi_hat * train[:-1]
sigma_hat = np.sqrt(np.mean(resid**2))
last_seen = np.empty_like(test)
last_seen[0] = train[-1]
last_seen[1:] = test[:-1]
ar_pred = phi_hat * last_seen  # forecast 1 step ahead with updating
rmse_ar = np.sqrt(np.mean((test - ar_pred)**2))
rmse_persist = np.sqrt(np.mean((test - last_seen)**2))
rmse_zero = np.sqrt(np.mean(test**2))
print("estimated AR parameter:", phi_hat)
print("residual standard deviation:", sigma_hat)
print("RMSE AR, persistence, zero:", rmse_ar, rmse_persist, rmse_zero)
print(acorr_ljungbox(resid, lags=[5, 10], return_df=True))

# Forecast from the last available observation at time n-1,
# with parameters fixed. Parameter-estimation uncertainty omitted.
h = 3
mean_h = phi_hat**h * x[-1]
sd_h = sigma_hat * np.sqrt(sum(phi_hat**(2*j) for j in range(h)))
ci_h = (mean_h - norm.ppf(0.975)*sd_h,
        mean_h + norm.ppf(0.975)*sd_h)
print("3-step forecast and approximate interval:", mean_h, ci_h)
\end{lstlisting}

\paragraph{Praktyczna interpretacja.}
Ocenę predykcji wykonujemy wyłącznie na obserwacjach
późniejszych niż trening. Warto porównać model AR(1)
z prognozą stałej średniej, prognozą ostatniej wartości,
a dla danych sezonowych także z prognozą sezonowo naiwną.
Dla ziarna generatora użytego w programie otrzymujemy
$\widehat\phi\approx0{,}731$ i RMSE prognozy AR(1)
około $1{,}014$. Wynik nie gwarantuje przewagi
na innych zbiorach danych.
W programie użyto \emph{aktualizacji} przy prognozie
jednokrokowej: do przewidzenia $X_t$ wolno użyć już
zaobserwowanego $X_{t-1}$, ale nie wartości $X_t$
ani późniejszych. Jeśli wymagamy jednorazowej prognozy
całych $40$ przyszłych okresów, trzeba zastosować
prognozę wielokrokową z \emph{jednego} momentu startowego,
która jest innym zadaniem i zwykle ma większy błąd.

\section{Zależne zdarzenia: model Markowa jako alternatywa dla iid}
\begin{przyklad}[Awaria w kolejnych dniach i trwałe stany]
Niech $S_t\in\{0,1\}$ oznacza stan maszyny w dniu $t$:
$0$ to stan prawidłowy, $1$ to stan alarmowy. Załóżmy,
że $(S_t)$ jest jednorodnym łańcuchem Markowa o macierzy
przejścia
$P=\begin{pmatrix}1-a&a\\ b&1-b\end{pmatrix}$,
gdzie $0<a,b<1$. Prawdopodobieństwo alarmu w stanie
stacjonarnym wynosi $\pi_1=a/(a+b)$, lecz kolejne
alarmy nie są niezależne, chyba że $a+b=1$.
Warunkowa wiarygodność obserwowanej ścieżki
$s_1,\ldots,s_n$ jest proporcjonalna do
$(1-a)^{N_{00}}a^{N_{01}}b^{N_{10}}(1-b)^{N_{11}}$,
gdzie $N_{ij}=\#\{t\in\{1,\ldots,n-1\}:s_t=i,s_{t+1}=j\}$.
W konsekwencji, jeśli odpowiednie mianowniki są dodatnie,
warunkowe MLE wynoszą
$\widehat a=N_{01}/(N_{00}+N_{01})$ i
$\widehat b=N_{10}/(N_{10}+N_{11})$.
Wyprowadzenie wynika z rozdzielenia log-wiarygodności
na dwa składniki dwumianowe; gdy obserwujemy zerową
liczbę przejść danego typu, trzeba uwzględnić brzeg
przestrzeni parametrów. Próba zignorowania kolejności
i dopasowania modelu niezależnych wskaźników Bernoulliego
prowadziłaby do innego modelu niepewności.
\end{przyklad}

\begin{stwierdzenie}[Korelacja w dwustanowym łańcuchu]
\label{prop:markov-corr}
W powyższym łańcuchu przy rozkładzie początkowym
stacjonarnym zachodzi
$\rho(h)=(1-a-b)^h$ dla $h\ge0$.
\end{stwierdzenie}
\begin{proof}
Oznaczmy $\lambda=1-a-b$.
Dla $s\in\{0,1\}$ mamy
$\E[S_{t+1}\mid S_t=s]=a+(1-a-b)s
=\pi_1+\lambda(s-\pi_1)$.
Iteracja własności Markowa daje
$\E[S_{t+h}-\pi_1\mid S_t]=\lambda^h(S_t-\pi_1)$.
Stąd
$\Cov(S_t,S_{t+h})=
\E[(S_t-\pi_1)\E[S_{t+h}-\pi_1\mid S_t]]
=\lambda^h\Var(S_t)$, co kończy dowód.
\end{proof}

\section{Dobór narzędzia i błędy praktyczne}
\begin{center}\small
\begin{tabular}{>{\raggedright\arraybackslash}p{.25\textwidth}>{\raggedright\arraybackslash}p{.32\textwidth}>{\raggedright\arraybackslash}p{.30\textwidth}}
\toprule
\textbf{Sytuacja} & \textbf{Metoda do rozważenia} & \textbf{Co zweryfikować}\\
\midrule
Jedna trajektoria o stałej średniej, krótka pamięć & AR/MA/ARMA & Stacjonarność, rzędy i ACF reszt.\\
Trend i sezonowość & Regresja z czasem, sezonowe ARIMA & Brak przecieku i stabilność reżimu.\\
Średnia danych zależnych & Długookresowy SE, HAC, model procesu & Sumowalność korelacji i CLT.\\
Dwie wartości stanu, istotna kolejność & Łańcuch Markowa & Jednorodność przejść, liczba przejść.\\
Zliczenia zdarzeń w kolejnych okresach & Poisson tylko przy uzasadnionej niezależności; modele zależnych zliczeń w pozostałych sytuacjach & Sezonowość, nadmierna dyspersja, korelacja.\\
Prognoza w czasie & Rolling origin, odłożona końcówka szeregu & Wyłącznie historyczne cechy, właściwy horyzont.\\
\bottomrule
\end{tabular}
\end{center}

\begin{zadanie}[Dlaczego okna mogą być skorelowane mimo procesu Poissona?]
Niech $(N_t)_{t\ge0}$ będzie jednorodnym procesem
Poissona o intensywności $\lambda>0$. Definiujemy
$A=N_2-N_0$ i $B=N_3-N_1$. Czy $A$ i $B$ są niezależne?
\par\emph{Rozwiązanie.} Nie. Zapiszmy
$A=(N_1-N_0)+(N_2-N_1)$ oraz
$B=(N_2-N_1)+(N_3-N_2)$.
Rozłączne przyrosty są niezależne, a wspólny składnik
$N_2-N_1\sim\Pois(\lambda)$ ma wariancję $\lambda$.
Stąd $\Cov(A,B)=\lambda>0$.
Model niezależnych \emph{nachodzących na siebie} okien
zliczeń byłby błędny, chociaż bazowy proces ma
niezależne przyrosty na rozłącznych przedziałach.
\end{zadanie}

\begin{zadanie}[Dobór błędu standardowego]
Dla $n=200$ kolejnych obserwacji stabilnego procesu
oszacowano autokorelację jako w przybliżeniu
$\rho(h)=0{,}6^{|h|}$. Analityk obliczył błąd
standardowy średniej $s/\sqrt n$. Jaką poprawkę
motywuje model AR(1)?
\par\emph{Rozwiązanie.} W przybliżeniu
$\tau^2/\gamma(0)=(1+0{,}6)/(1-0{,}6)=4$.
Błąd standardowy średniej jest około dwa razy
większy od klasycznego, o ile model i przybliżenie
asymptotyczne są odpowiednie. W analizie danych należy
oszacować parametr, wariancję długookresową i ocenić
niepewność tych oszacowań; nie wolno stosować
tego mnożnika bez sprawdzenia struktury zależności.
\end{zadanie}


% ================================================================
% ROZDZIAL 17 -- WNIOSKOWANIE BAYESOWSKIE
% ================================================================
